% Part: turing-machines
% Chapter: machines-computations
% Section: variants

\documentclass[../../../include/open-logic-section]{subfiles}

\begin{document}

\olfileid[hi]{tur}{mac}{var}
\olsection{ट्यूरिंग मशीनों के रूपांतर}

वास्तव में ट्यूरिंग मशीनों को परिभाषित करने की अनेक संभव रीतियाँ हैं;
हमारी रीति उनमें से केवल एक है। कुछ दृष्टियों से हमारी परिभाषा अन्य
परिभाषाओं की अपेक्षा अधिक उदार है। हम मनमानी परिमित वर्णमालाओं की अनुमति
देते हैं; कोई अधिक प्रतिबंधित परिभाषा पट्टी पर केवल दो प्रतीकों
$\TMstroke$ और~$\TMblank$ की अनुमति दे सकती है। हम मशीन को पट्टी पर
प्रतीक लिखने और उसी समय चलने देते हैं; अन्य परिभाषाएँ लिखने अथवा चलने में
से केवल एक की अनुमति देती हैं। हम पट्टी के शीर्ष को चलाए बिना लिखने की
संभावना स्वीकार करते हैं; अन्य परिभाषाएँ $\TMstay$ ``निर्देश'' को छोड़
देती हैं। दूसरी दृष्टियों से हमारी परिभाषा अधिक प्रतिबंधित है। हमने माना
कि पट्टी केवल एक दिशा में अनंत है; अन्य परिभाषाएँ पट्टी को बाईं और दाईं
दोनों ओर अनंत होने देती हैं। वास्तव में, कितनी भी अलग पट्टियों, या खानों
के अनंत जालक तक की अनुमति दी जा सकती है। हम ट्यूरिंग मशीन के
निर्देश-समुच्चय को संक्रमण फलन द्वारा निरूपित करते हैं; अन्य परिभाषाएँ
संक्रमण संबंध का उपयोग करती हैं, जिसमें किसी दी हुई स्थिति में मशीन के
लिए एक से अधिक संभव निर्देश होते हैं।

परिभाषा का यह अंतिम शिथिलीकरण विशेष रूप से रोचक है। हमारी परिभाषा में,
जब मशीन अवस्था~$q$ में प्रतीक~$\sigma$ पढ़ रही होती है, तब
$\delta(q, \sigma)$ नया प्रतीक, अवस्था और पट्टी के शीर्ष की स्थिति
निर्धारित करता है। किंतु यदि निर्देश-समुच्चय को वर्तमान अवस्था-प्रतीक
युग्मों $\tuple{q,
  \sigma}$ और नए अवस्था-प्रतीक-दिशा त्रिकों $\tuple{q', \sigma',
  D}$ के बीच संबंध होने दें, तो ट्यूरिंग मशीन की
क्रिया अद्वितीय रूप से निर्धारित न हो---निर्देश संबंध में
$\tuple{q,
  \sigma, q', \sigma', D}$ और
$\tuple{q, \sigma, q'', \sigma'', D'}$ दोनों हो सकते हैं। इस स्थिति में
हमारे पास एक \emph{अनियतात्मक} ट्यूरिंग मशीन होती है। संगणनात्मक जटिलता
सिद्धांत में ऐसी मशीनों की महत्त्वपूर्ण भूमिका है।

ट्यूरिंग मशीन कब रुकती है, इसके लिए भी भिन्न परिपाटियाँ हैं: हम कहते
हैं कि संक्रमण फलन अपरिभाषित होने पर वह रुकती है; अन्य परिभाषाएँ मशीन
का किसी विशेष निर्दिष्ट विराम-अवस्था में होना आवश्यक मानती हैं।
\olref[hal]{sec} में हम समझा चुके हैं कि निर्दिष्ट विराम-अवस्था की
आवश्यकता ऐसा प्रतिबंध क्यों नहीं है जो ट्यूरिंग मशीनों की संगणन-क्षमता
को प्रभावित करे। चूँकि हमारी ट्यूरिंग मशीनों की पट्टियाँ केवल एक दिशा
में अनंत हैं, इसलिए ऐसी स्थितियाँ आती हैं जहाँ ट्यूरिंग मशीन किसी
निर्देश को ठीक से निष्पादित नहीं कर सकती: जब वह सबसे बायाँ खाना पढ़ रही
हो और उसे बाईं ओर जाना हो। हमारी परिभाषा के अनुसार वह ``गिरने'' के
स्थान पर वहीं रहती है; किंतु हम इसे ऐसी परिभाषा दे सकते थे कि ऐसा होने
पर वह रुक जाए। यह परिभाषा भी समतुल्य है: खाना~$0$ से बाईं ओर जाने का
प्रयास करने पर रुकने वाली ट्यूरिंग मशीन के व्यवहार का अनुकरण प्रत्येक
संक्रमण $\delta(q,\TMendtape) =
\tuple{q',\sigma,\TMleft}$ को मिटाकर कर
सकते हैं---तब~$\TMendtape$ पर बाईं ओर जाने का प्रयास करने के स्थान पर
मशीन रुक जाती है।\footnote{यह \emph{पूरी तरह} काम नहीं करता, क्योंकि
हमें $\TMendtape$ उन खानों पर भी लिखने और पढ़ने से कुछ नहीं रोकता जो
खाना~$0$ नहीं हैं (देखें \olref[una]{ex:mover})। ऐसे प्रयोजन के लिए
इसके स्थान पर उपयोग करने हेतु दूसरा~$\TMendtape'$ प्रतीक जोड़कर इसे हल
कर सकते हैं।}

संख्याओं के निरूपण की भी भिन्न रीतियाँ हैं (और इसलिए ट्यूरिंग मशीन द्वारा
संगणित आगत-निर्गत फलन की भी): हम एकाधारी निरूपण का उपयोग करते हैं, किंतु
द्विआधारी निरूपण भी उपयोग कर सकते हैं। इसके लिए $\TMblank$ और
~$\TMendtape$ के अतिरिक्त दो प्रतीक चाहिए।

अब एक रोचक तथ्य देखें: कौन-से फलन ट्यूरिंग-संगणनीय हैं, इस पर इनमें से
किसी रूपांतर का प्रभाव नहीं पड़ता। \emph{यदि कोई फलन एक परिभाषा के
अनुसार ट्यूरिंग-संगणनीय है, तो वह उन सभी के अनुसार ट्यूरिंग-संगणनीय है।}

हम इसके सत्यापन के विस्तार में नहीं जाएँगे। केवल एक उदाहरण देखें:
दोनों दिशाओं में अनंत पट्टी अथवा अनेक पट्टियों की अनुमति देने से कोई
अतिरिक्त संगणन-शक्ति नहीं मिलती। मोटे तौर पर कारण यह है कि एकल एकदिश
अनंत पट्टी वाली ट्यूरिंग मशीन अनेक पट्टियों अथवा द्विदिश अनंत पट्टियों
का अनुकरण कर सकती है। उदाहरणार्थ, अतिरिक्त अवस्थाओं और निर्देशों का
उपयोग करके अनेक पट्टियों या द्विदिश अनंत पट्टी वाली मशीन के प्रोग्राम का
``अनुवाद'' एकल एकदिश अनंत पट्टी वाली मशीन में कर सकते हैं। अनूदित मशीन
सम खानों को पट्टी~$1$ के खानों (या द्विदिश अनंत पट्टी के ``धन'' खानों)
के लिए और विषम खानों को पट्टी~$2$ के खानों (या ``ऋण'' खानों) के लिए
उपयोग कर सकती है।

\end{document}
